\section{Proofs for the variants: great hyperspheres, points anywhere, flats, convex bodies}
\label{app:D}

% ---------------------------------------------------------------------------
\subsection{Great hyperspheres through \texorpdfstring{$d$}{d} points}
\label{sec:D-spheres}

\begin{definition}\leavevmode
\begin{itemize}
  \item $\Sph^d\subset\R^{d+1}$ carries the angular metric. A linearly independent $d$-subset $F$ of a
    finite $P\subset\Sph^d$ spans the hyperplane $\spn(F)$ through $0$, and gives the great hypersphere
    $\Sph^d\cap\spn(F)$, at angular distance $\arcsin\abs{\ip{x}{\nu}}$ from $x\in\Sph^d$, $\nu$ a unit normal.
  \item For unit vectors $\nu_1,\dots,\nu_k$, $\rho(\nu)=\arcsin\max_{x\in\Sph^d}\min_j\abs{\ip{x}{\nu_j}}$ (the max is
    attained, and the signs of the $\nu_j$ do not matter). $\rho_S(P)$ is $\rho$ of the normals of the great
    hyperspheres of $P$ ($+\infty$ if there are none), and $\lambda_S^{(d)}(n)=\inf\rho_S(P)$ over $n$-point
    $P\subset\Sph^d$.
  \item $p_d(n)$ is as in Theorem~B, and $K=\binom nd$.
\end{itemize}
\end{definition}

\begin{proposition}\label{D:spheres}
For $d\ge2$ and $n>d$, $\pi/(2K)\le\lambda_S^{(d)}(n)\le\pi/(2p_d(n))$. (The Euclidean distance from $x\in\Sph^d$ to $\spn(F)$ is $\abs{\ip{x}{\nu}}$, the sine of the angular one; as
$\rho\le\pi/2$ and $\sin$ increases on $[0,\pi/2]$, the Euclidean covering radius of $P$ is $\sin\rho_S(P)$, and its
infimum lies between $\sin(\pi/(2K))$ and $\sin(\pi/(2p_d(n)))$.)
\end{proposition}

\begin{proof}
\textbf{Lower bound.} If $\rho=\rho_S(P)<\infty$, the closed angular $\rho$-neighbourhoods of the at most $K$
great hyperspheres are spherical segments of (angular) width $2\rho$ covering $\Sph^d$. By Glazyrin--Karasev--Polyanskii
(arXiv:2112.05382, Definition~1.14 and Corollary~1.15: if $\Sph^{m-1}$, $m\ge2$, is covered by spherical
segments, their angular widths sum to at least $\pi$), applied with $m=d+1$, $2\rho K\ge\pi$.

\textbf{Upper bound.} Let $N=p_d(n)$. The construction of Theorem~B (via Lemma~2) gives $n$-point
configurations $Q_s\subset\R^d$ ($s>0$) and a family $\mathcal F$ of $d$-subsets, affinely independent in
every $Q_s$, such that as $s\to0$ the hyperplane $\{y:\ip{\mu_F}{y}=c_F\}$ ($\abs{\mu_F}=1$) of each
$F\in\mathcal F$ converges to a hyperplane through $0$ containing a fixed $(d-2)$-subspace $L$: after a choice of
the sign of $(\mu_F,c_F)$, $\mu_F\to\mu_F^0$ and $c_F\to0$; and the limit normals $\mu_F^0$ lie in the plane $L^\perp$ with directions mod $\pi$
having largest gap $G\le\pi/N$. Lift $Q_s$ to $\Sph^d$ by $y\mapsto(y,1)/\abs{(y,1)}$; a $d$-subset is linearly
independent after the lift iff it was affinely independent, and the lift of $F$ spans the hyperplane with
unit normal $(\mu_F,-c_F)/\sqrt{1+c_F^2}\to(\mu_F^0,0)$.

For these limit normals and $x\in\Sph^d$, let $x_2$ be the component of $x$ in $L^\perp\times\{0\}$. Then
$\min_F\abs{\ip{x}{(\mu_F^0,0)}}=\min_F\abs{\ip{x_2}{\mu_F^0}}$. This is $0$ if $x_2=0$; otherwise it is
$\abs{x_2}\min_F\abs{\cos\angle(x_2,\mu_F^0)}=\abs{x_2}\min_F\sin\angle(x_2,\ell_F)\le\sin(G/2)$, where $\ell_F\subset L^\perp$ is
the line orthogonal to $\mu_F^0$: the lines $\ell_F$ also have largest gap $G$ mod $\pi$, so one is within $G/2$
of $x_2$, and $G/2\le\pi/2$ (gaps mod $\pi$ are at most $\pi$, $\mathcal F\ne\emptyset$); also $\abs{x_2}\le1$. Equality holds
for the unit $x=x_2$ at the middle of a largest gap. So the limit family has $\rho=G/2\le\pi/(2N)$. Finally,
$\sin\rho(\nu)$ is $1$-Lipschitz for the distance $\max_j\min_\pm\abs{\nu_j\mp\nu'_j}$ between families indexed alike
(since $\abs{\abs{\ip{x}{a}}-\abs{\ip{x}{b}}}\le\abs{a-b}$ for unit $x$) and $\arcsin$ is uniformly continuous,
so $\rho$ of the normals of the lifts of $F\in\mathcal F$ tends to $G/2$; the other great hyperspheres only
lower $\rho_S$. Hence $\lambda_S^{(d)}(n)\le\pi/(2N)$.
\end{proof}

% ---------------------------------------------------------------------------
\subsection{Points anywhere: \texorpdfstring{$\mu_d(n)=\Theta(n^{-d})$}{mu\_d(n) = Theta(n\textasciicircum -d)} in every dimension, and \texorpdfstring{$\mu_d(d+1)=1/(d+1)$}{mu\_d(d+1) = 1/(d+1)}}
\label{sec:D-anywhere}

This proves the bounds for points anywhere in every dimension $d\ge2$.

\begin{definition}\leavevmode
\begin{itemize}
  \item $\ball^d$ is the closed unit ball. For a finite $P\subset\R^d$ (anywhere, not only in
    $\ball^d$), $\mathcal{H}_P$ is the set of hyperplanes through $d$ affinely independent points
    of $P$, and
    \[
      \rho(P)=\max_{x\in\ball^d}\ \min_{H\in\mathcal{H}_P}\dist(x,H),
    \]
    with $\rho=\infty$ if $\mathcal{H}_P=\emptyset$.
  \item $\mu_d(n)=\inf\{\rho(P):P\subset\R^d,\ \abs{P}=n\}$, for $d\ge2$ and $n\ge d$.
  \item For a finite set $S$ of positive integers, $N_d(S)$ is the length of the longest run of
    consecutive integers all of which are sums of $d$ distinct elements of $S$.
    $N_d(n)=\max_{\abs{S}=n}N_d(S)$.
\end{itemize}
\end{definition}

\begin{proposition}\label{D:anywhere}
For $d\ge2$ and $n\ge d$:
\begin{enumerate}[label=(\roman*)]
  \item $1/\binom{n}{d}\le\mu_d(n)\le1/N_d(n)$.
  \item $\mu_d(d+1)=1/(d+1)$.
  \item $N_d(dq)\ge q^d$ for $q\ge2$. Since $\mu_d$ is non-increasing in $n$,
    \[
      d!\le\liminf n^d\mu_d(n)\le\limsup n^d\mu_d(n)\le d^d.
    \]
\end{enumerate}
\end{proposition}

\begin{proof}
\textbf{Lower bound.} Write the $k'\le\binom{n}{d}$ hyperplanes of $\mathcal{H}_P$ as
$\{x:\ip{x}{u_j}=m_j\}$ with $\abs{u_j}=1$. Lemma~1 (Bang's sign argument: if the closed
$\eps$-slabs around these hyperplanes cover $\ball^d$, then
$\eps\cdot\max_{\abs v=1}\sum_j\abs{\ip{u_j}{v}}\ge1$) does not need the points to lie in the ball.
The max defining $\rho(P)$ is attained ($\ball^d$ is compact), so the closed $\rho(P)$-slabs cover
$\ball^d$. Since $\abs{\ip{u_j}{v}}\le1$ for unit $v$,
$\rho k'\ge\rho\cdot\max_{\abs v=1}\sum_j\abs{\ip{u_j}{v}}\ge1$. So $\rho(P)\ge1/\binom{n}{d}$.

\begin{lemma}[Interpolation identity]\label{D:interp}
Let $t_1,\dots,t_d$ be distinct and nonzero, and let $p$ be the polynomial of degree $\le d-1$ with
$p(t_i)=1/t_i$. Then $p(0)=\sum_i1/t_i$.
\end{lemma}

\begin{proof}
The polynomial $1-t\,p(t)$ has degree $\le d$ and vanishes at every $t_i$, so
$1-t\,p(t)=c\prod_i(t-t_i)$. At $t=0$ this gives $c=1/\prod(-t_i)$. The coefficient of $t$ gives
$-p(0)=c\cdot(-1)^{d-1}e_{d-1}$, where $e_k$ are the elementary symmetric functions of the $t_i$. So
$p(0)=e_{d-1}/e_d=\sum1/t_i$.
\end{proof}

\textbf{Upper bound (i).} Let $S=\{u_1,\dots,u_n\}$ be distinct positive integers whose
$d$-distinct sums contain $c+1,\dots,c+N$. Put $t_i=1/u_i$, $\alpha=2/N$ and
$\beta=-1-\alpha c-\alpha/2$, so that $\alpha(c+j)+\beta=-1+(2j-1)/N$ for $j=1,\dots,N$. For $X>0$
let
\begin{equation*}
  p_i(X)=\bigl(Xt_i,\,Xt_i^2,\,\dots,\,Xt_i^{d-1},\,\alpha u_i+\beta\bigr)\in\R^d\qquad(i=1,\dots,n).
\end{equation*}

Write points of $\R^d$ as $(z,y)$ with $z\in\R^{d-1}$. Fix a $d$-set $F$, and let $p_F$ be the
interpolating polynomial of degree $\le d-1$ with $p_F(t_i)=\alpha/t_i+\beta$ for $i\in F$. Write
$p_F(t)=b_0+b_1t+\dots+b_{d-1}t^{d-1}$. Then the affine function
$g_F(z)=b_0+\sum_{k=1}^{d-1}(b_k/X)z_k$ satisfies $g_F(Xt_i,\dots,Xt_i^{d-1})=p_F(t_i)$, so its
graph contains $p_i(X)$ for $i\in F$.

The points $p_i(X)$, $i\in F$, are affinely independent. Their first $d-1$ coordinates are
$X\cdot(t_i,\dots,t_i^{d-1})$, which are affinely independent by Vandermonde. So the graph of $g_F$
is the hyperplane $H_F(X)$ through them, and every $d$-set is independent.

By the identity (Lemma~\ref{D:interp}), $b_0=p_F(0)=\alpha\sum_{i\in F}u_i+\beta$. The $b_k$ do not
depend on $X$. So as $X\to\infty$, $H_F(X)$ converges (unit normal and offset) to the horizontal
hyperplane $\{y=\alpha\sum_{i\in F}u_i+\beta\}$.

The covering radius of a fixed finite list of hyperplanes is continuous in their normals and
offsets: for $\abs{x}\le1$, $\abs{\dist(x,H)-\dist(x,H')}\le\abs{n-n'}+\abs{c-c'}$, and max-min
preserves this bound. So $\rho(P(X))$ tends to the covering radius of the limit list. That list
contains the horizontal hyperplanes at heights $-1+(2j-1)/N$, $j=1,\dots,N$, whose $1/N$-slabs cover
$\ball^d$. Hence $\mu_d(n)\le\lim\rho(P(X))\le1/N$.

\textbf{(ii)} Take $S=\{1,\dots,d+1\}$. Its $d$-distinct sums are $T-s$ for $s\in S$, where
$T=\sum S$. These are $d+1$ consecutive integers, so $N_d(d+1)\ge d+1$. Together with the lower
bound $1/\binom{d+1}{d}$, $\mu_d(d+1)=1/(d+1)$.

\textbf{(iii)} Let $L=q^d$ and
\[
  S=\bigcup_{j=0}^{d-1}\{aq^j+jL+1:0\le a<q\}:
\]
$d$ blocks of $q$ elements, $n=dq$. Elements within a block are distinct since $q^j\ge1$, and the
blocks are disjoint, since (using $(q-1)q^j+1\le q^d$) block $j$ lies in
$[jL+1,\,jL+1+(q-1)q^{d-1}]\subset[jL+1,\,(j+1)L]$. Choosing one element from each block gives
distinct elements with sum $\sum_ja_jq^j+\sum_j(jL+1)$. As $(a_0,\dots,a_{d-1})$ ranges over
$\{0,\dots,q-1\}^d$, $\sum a_jq^j$ takes all values $0,\dots,q^d-1$. So $N_d(dq)\ge q^d$ and
$\mu_d(dq)\le q^{-d}$.

Adding a point never removes a hyperplane, so $\mu_d$ is non-increasing in $n$. For
$dq\le n<d(q+1)$, $n^d\mu_d(n)\le(d(q+1))^d/q^d\to d^d$. The liminf bound is
$n^d/\binom{n}{d}\to d!$.
\end{proof}

\begin{lemma}[The plane]\label{D:N2}
$N_2(n)\ge\lfloor n^2/4\rfloor$ for $n\ge2$, so $\mu_2(n)\le1/\lfloor n^2/4\rfloor$. Moreover $\mu_2(3)=1/3$
and $\mu_2(4)=1/6$.
\end{lemma}

\begin{proof}
Let $a=\lfloor n/2\rfloor$, $A=\{1,\dots,a\}$ and $B=\{2a,3a,\dots,(n-a+1)a\}$: $n$ distinct positive
integers. The sums $i+ja$ ($1\le i\le a$, $2\le j\le n-a+1$) of distinct elements fill the interval
$[2a+1,(n-a+2)a]$, which contains $a(n-a)=\lfloor n^2/4\rfloor$ integers. With
Proposition~\ref{D:anywhere}(i) this gives $\mu_2(n)\le1/\lfloor n^2/4\rfloor$. $\mu_2(3)=1/3$ is
Proposition~\ref{D:anywhere}(ii) with $d=2$. The pair sums of
$\{1,2,3,5\}$ are $3,\dots,8$, so $1/6=1/\binom42\le\mu_2(4)\le1/N_2(4)\le1/6$.
\end{proof}

% ---------------------------------------------------------------------------
\subsection{Flats through \texorpdfstring{$j+1$}{j+1} points: \texorpdfstring{$\lambda_d^{(j)}(n)=\Theta(n^{-(j+1)/(d-j)})$}{lambda = Theta(n\textasciicircum(-(j+1)/(d-j)))} for every \texorpdfstring{$1\le j\le d-1$}{1 <= j <= d-1}}
\label{sec:D-flats}

For $j=1$ (lines) this section proves the tube bound $k\rho^{d-1}\ge1$ and, in $\ball^3$,
$\sqrt2\le\liminf n\lambda\le\limsup n\lambda\le6\sqrt2$. For $j=d-1$ (hyperplanes) it recovers the
order $n^{-d}$, with weaker constants than the main text. Floating-point exploration (not part of the
proofs): the accompanying repository
(\texttt{experiments/\allowbreak flats\_\allowbreak middle\_\allowbreak dims.py}).

\begin{definition}\leavevmode
\begin{itemize}
  \item $d\ge2$ and $1\le j\le d-1$. Put $m=d-j+1\ge2$. $\ball^d$ is the closed unit ball,
    $\Sph^{k-1}$ the unit sphere of $\R^k$, $\kappa_k=\pi^{k/2}/\Gamma(k/2+1)$ and
    $\abs{\Sph^{k-1}}=2\pi^{k/2}/\Gamma(k/2)$.
  \item A $j$-flat is an affine $j$-dimensional subspace. For $P\subset\ball^d$ with $n\ge j+1$
    points, $\mathcal{F}_P$ is the set of $j$-flats spanned by affinely independent
    $(j+1)$-subsets of $P$, so $\abs{\mathcal{F}_P}\le K=\binom{n}{j+1}$. Put
    \[
      \rho(P)=\max_{x\in\ball^d}\ \min_{F\in\mathcal{F}_P}\dist(x,F)\quad(=\infty\text{ if }\mathcal{F}_P=\emptyset),
    \]
    and $\lambda_d^{(j)}(n)=\inf\{\rho(P):P\subset\ball^d,\ \abs{P}=n\}$.
  \item For a finite set $W\subset\R^m\setminus\{0\}$,
    $\delta(W)=\max_{x\in\Sph^{m-1}}\min_{w\in W}\angle(x,\pm w)\in[0,\pi/2]$.
\end{itemize}
\end{definition}

\begin{proposition}[Lower bound]\label{D:M1}
If the $\rho$-neighbourhoods of $K$ $j$-flats cover $\ball^d$, then $K\rho^{d-j}\ge c(d,j)$, where
\begin{equation*}
  c(d,1)=1,\qquad c(d,j)=\frac{\Gamma(j/2)\,\Gamma((d-j)/2+1)}{\Gamma(d/2)}\quad(j\ge2).
\end{equation*}
Hence $\binom{n}{j+1}\cdot\lambda_d^{(j)}(n)^{d-j}\ge c(d,j)$.
\end{proposition}

For $j=1$ this is the tube bound $k\rho^{d-1}\ge1$ for $k$ lines. For $(d,j)=(3,2)$ it is the plank bound $K\rho\ge1$. For
$j\ge2$ it beats the volume bound $\kappa_d/(\kappa_j\kappa_{d-j})$ by the factor $d/j$.

\begin{proof}
For $0\le a<1$ let $f_a(x)=(1-\abs{x}^2)^{-a}$ on the open ball.
\begin{enumerate}[label=(\alph*)]
  \item \textbf{A flat.} Let $F$ be a $j$-flat at distance $h<1$ from $0$, and $R=\sqrt{1-h^2}$.
    Then
    \[
      \int_{F\cap\ball^d}f_a=\abs{\Sph^{j-1}}\int_0^Rr^{j-1}(R^2-r^2)^{-a}\,dr
      =R^{j-2a}\cdot\tfrac12\abs{\Sph^{j-1}}\,B(j/2,1-a).
    \]
    If $a\le j/2$, this is at most $I_j(a):=\frac12\abs{\Sph^{j-1}}\,B(j/2,1-a)$, its value at
    $h=0$. Flats with $h\ge1$ miss the open ball.
  \item \textbf{The ball.} $\int_{\ball^d}f_a=\frac12\abs{\Sph^{d-1}}\,B(d/2,1-a)$.
  \item \textbf{A neighbourhood.} The $\rho$-neighbourhood of $F$ is the union of the parallel flats
    $F+w$, for $w$ in the $(d-j)$-dimensional orthogonal complement of the direction of $F$ with
    $\abs{w}\le\rho$. By Fubini and (a), its $f_a$-mass is at most $\kappa_{d-j}\rho^{d-j}I_j(a)$.
  \item \textbf{Covering.} $K\kappa_{d-j}\rho^{d-j}I_j(a)\ge\frac12\abs{\Sph^{d-1}}\,B(d/2,1-a)$.
    Simplifying with $B(x,y)=\Gamma(x)\Gamma(y)/\allowbreak\Gamma(x+y)$:
    \begin{equation*}
      K\rho^{d-j}\ge c_a:=\frac{\Gamma(j/2+1-a)\,\Gamma((d-j)/2+1)}{\Gamma(d/2+1-a)},
    \end{equation*}
    for every $a<1$ with $a\le j/2$.
    For $j=1$ take $a=\frac12$: $c_{1/2}=1$. For $j\ge2$ let $a\uparrow1$: $c_a\to c(d,j)$. The bound
    holds for every $a<1$, so it holds in the limit. With $a=0$,
    $c_0=\kappa_d/(\kappa_j\kappa_{d-j})$ is the volume bound, and $c(d,j)/c_0=d/j$ for
    $j\ge2$.\qedhere
\end{enumerate}
\end{proof}

\begin{lemma}[Flat clusters through a $(j-1)$-space]\label{D:M2}
Write $\R^d=\R^{j-1}\times\R^m$. Let $\tau_0<\dots<\tau_j$ be reals,
$t_r=\eta\cdot(\tau_r,\tau_r^2,\dots,\tau_r^{j-1})\in\R^{j-1}$ (empty if $j=1$) with $\eta>0$ (a
scale),
and $c_r=1/\prod_{k\ne r}(\tau_r-\tau_k)$. Then $\sum_rc_r=0$ and $\sum_rc_rt_r=0$, and $c$ spans
the affine dependencies of $t_0,\dots,t_j$. For $a_r\in\R^m$ with $w=\sum_rc_ra_r\ne0$, and $s>0$,
the points $q_r=(t_r,sa_r)$ are affinely independent. As $s\to0$, their $j$-flat converges (in the
sense that $x\mapsto\dist(x,\cdot)$ converges uniformly on $\ball^d$) to the linear subspace
$L_w=(\R^{j-1}\times\{0\})\oplus\R(0,w)$.
\end{lemma}

\begin{proof}
\emph{The dependency.} The $j$-th divided difference $\sum_rc_rp(\tau_r)$ of a polynomial of degree
$<j$ vanishes. So $\sum_rc_rp(\tau_r)=0$ for $p=1,\tau,\dots,\tau^{j-1}$, which gives $\sum c_r=0$
and $\sum c_rt_r=0$. The $t_r$ affinely span $\R^{j-1}$: any $j$ of them have Vandermonde determinant
$\ne0$. So the dependency space is one-dimensional and spanned by $c$. The rest is Lemma~2 for
$j$-flats (Section~\ref{sec:B-lemma2}, parts (i) and (ii)), applied to these $t_r$, $c$ and $a_r$.
\end{proof}

For linear $L_w$ and any $x$, $\dist(x,L_w)=\abs{x_\perp}\sin\angle(x_\perp,\R w)$, where $x_\perp$
is the $\R^m$-component (read as $0$ when $x_\perp=0$). Hence
\begin{equation*}
  \max_{x\in\ball^d}\ \min_{w\in W}\dist(x,L_w)=\sin\delta(W).
\end{equation*}

\begin{proposition}[Construction]\label{D:M3}
Let $M\ge2$ and $N=M^{j+1}$, and assume $N-1\ge(2/\pi)\sqrt{m-1}$, so that
$\theta:=\sqrt{m-1}/(N-1)\le\pi/2$. Then
\begin{equation*}
  \lambda_d^{(j)}\bigl(m(j+1)M^{m-1}\bigr)\le\sin\theta.
\end{equation*}
Adding a point keeps every flat already spanned, so $\lambda$ is non-increasing in $n$; hence
\begin{equation*}
  \limsup_n n^{(j+1)/(d-j)}\lambda_d^{(j)}(n)\le\sqrt{d-j}\cdot\bigl((d-j+1)(j+1)\bigr)^{(j+1)/(d-j)}.
\end{equation*}
\end{proposition}

For $j=1$, $d=3$ this is $6\sqrt2$. For $j=d-1$ it is $(2d)^d$.

\begin{proof}
\emph{Construction (one block).} Take $\tau_r=r$ ($r=0,\dots,j$), and $t_r$, $c_r$ as in
Lemma~\ref{D:M2}. Let $G=\{0,\dots,M-1\}^{m-1}$, $h=(N-1)/2$ and $e=(1,\dots,1)\in\R^{m-1}$. Group
$r$ consists of the points $(t_r,s\cdot b/c_r)$, where $b$ runs over
\begin{itemize}
  \item $r=0$: $B_0=\{(g-h\cdot e,\,h):g\in G\}\subset\R^{m-1}\times\R$;
  \item $r\ge1$: $B_r=\{(M^rg,\,0):g\in G\}$.
\end{itemize}
Each group has $M^{m-1}$ points, distinct because the $b$'s are distinct. Groups are distinct
because their $t_r$ differ. For $j=1$ there is no $t$; then $c_0=-1$ and $c_1=1$, group $0$ has last
coordinate $-s\cdot h\ne0$ and group $1$ has last coordinate $0$, so the groups are distinct too.

\emph{Transversal sets.} Take one point from each group, $F=\{(t_r,sb_r/c_r)\}$. Then
\[
  w_F=\sum_rc_r(b_r/c_r)=\sum_rb_r=(g_0+Mg_1+\dots+M^jg_j-h\cdot e,\,h).
\]
By base-$M$ digits, $g_0+\dots+M^jg_j$ runs over $\{0,\dots,N-1\}^{m-1}$. So the set $W$ of all
$w_F$ is $\{(x,h):x\in\{-h,-h+1,\dots,h\}^{m-1}\}$, and $w_F\ne0$. By Lemma~\ref{D:M2}, these flats
converge to $L_{w_F}$.

\emph{Covering one chart.} Let $x\in\Sph^{m-1}$ with $\abs{x_i}\le x_m$ for $i<m$. Then $x$ is a
positive multiple of $(p,1)$ with $p\in[-1,1]^{m-1}$. Writing $w=(z,h)$, the points
$w/h=(z/h,1)$ have $z/h$ on a grid of spacing $1/h$ covering $[-1,1]^{m-1}$, so some has
$\abs{p-z/h}\le\sqrt{m-1}/(2h)=\sqrt{m-1}/(N-1)$. By the gnomonic inequality
$\angle((p,1),(p',1))\le\abs{p-p'}$ for $p,p'\in\R^{m-1}$, the angle is at most $\sqrt{m-1}/(N-1)$.
(Proof: the path $P(s)=(p+s(p'-p),1)$, $s\in[0,1]$, has $\abs{P(s)}\ge1$; the angle is at most the
length of the spherical path $P/\abs P$, and $\abs{(P/\abs P)'}\le\abs{P'}/\abs P$, so the length is at
most $\int_0^1\abs{P'(s)}/\abs{P(s)}\,ds\le\abs{p'-p}$.)

\emph{All directions: $m$ blocks.} For $\sigma$ the cyclic permutation of the $m$ coordinates of
$\R^m$, block $i$ ($i=0,\dots,m-1$) applies $\sigma^i$ to the $\R^m$-components of block $0$, and
adds $s\cdot v_i$ to them, with $v_i\in\R^m$ chosen so that all points of all blocks are distinct
(this excludes finitely many affine conditions, so generic $v_i$ work). The translation does not
change $w_F$, since $\sum c_r=0$. Every $x\in\Sph^{m-1}$ is $\pm\sigma^i$ of a vector in the chart
above for some $i$. So $\delta(W_{\mathrm{all}})\le\sqrt{m-1}/(N-1)$, with $n=m(j+1)M^{m-1}$ points
in total.

\emph{Conclusion.} Take $\eta$ small enough that $\abs{t_r}<\frac12$, and $s$ small enough that all
points lie in $\ball^d$. Let $\mathcal{F}_s$ be the family of $j$-flats spanned by the transversal
sets of all blocks. Then $\rho(P_s)\le\max_x\min_{F\in\mathcal{F}_s}\dist(x,F)$, because extra flats
only lower $\rho$. By Lemma~\ref{D:M2} and uniform convergence, the right side tends to
$\sin\delta(W_{\mathrm{all}})\le\sin\theta$, using $\theta\le\pi/2$ (the hypothesis) and $\sin$
increasing on $[0,\pi/2]$. Without the hypothesis the bound can fail: for $(d,j,M)=(90,1,2)$,
$\theta>\pi$ and $\sin\theta<0$. For the limsup, put
$M=\bigl\lfloor(n/(m(j+1)))^{1/(m-1)}\bigr\rfloor$, which satisfies the hypothesis for large $n$, and
use monotonicity in $n$.
\end{proof}

\begin{corollary}\label{D:M4}
For all $d\ge2$ and $1\le j\le d-1$, $\lambda_d^{(j)}(n)=\Theta(n^{-(j+1)/(d-j)})$, with
\begin{align*}
  \bigl((j+1)!\cdot c(d,j)\bigr)^{1/(d-j)}&\le\liminf n^{(j+1)/(d-j)}\lambda\\
  &\le\limsup n^{(j+1)/(d-j)}\lambda\le\sqrt{d-j}\cdot\bigl((d-j+1)(j+1)\bigr)^{(j+1)/(d-j)}.
\end{align*}
\end{corollary}

The lower constant comes from Proposition~\ref{D:M1} with $\binom{n}{j+1}\sim n^{j+1}/(j+1)!$. For
example, for planes through triples in $\ball^4$ ($d=4$, $j=2$):
$\sqrt6\approx2.45\le\liminf n^{3/2}\lambda\le\limsup n^{3/2}\lambda\le27\sqrt2\approx38.2$.

\begin{remark}[Not part of the claims]\leavevmode
\begin{itemize}
  \item The construction is the direct analogue of Lemma~2 and Theorem~B: the $(j-1)$-space
    $\R^{j-1}\times\{0\}$ plays the role of $\R^{d-2}\times\{0\}$. Directions are sums of one vector
    from each group (here exactly, because $\tau$ is constant on each group), and digit grids make
    the sums a complete grid. For $j=d-1$ the additive bases of the main text give much better
    constants.
  \item The lower-bound measure $f_a$ concentrates on the sphere as $a\to1$. So for $j\ge2$,
    $c(d,j)$ is the bound obtained by covering $\Sph^{d-1}$ with neighbourhoods of the flats'
    $(j-1)$-subspheres.
  \item Floating point, not part of the proof (the accompanying repository
    (\texttt{experiments/\allowbreak flats\_\allowbreak middle\_\allowbreak dims.py})): the flats through transversal sets match $L_w$
    to $10^{-3}$ at $s=10^{-6}$. The measured one-chart $\delta$ is below the bound for
    $(d,j,M)=(4,2,3)$, $(5,2,2)$, $(5,3,2)$ and $(4,1,4)$.
\end{itemize}
\end{remark}

\begin{remark}[Lines through the centre of $\ball^3$]
If $k$ lines through $0$ have covering radius $\rho=\max_{x\in\ball^3}\min_\ell\dist(x,\ell)<1$, each $\rho$-tube meets $\Sph^2$ in two antipodal
caps of area $2\pi\bigl(1-\sqrt{1-\rho^2}\bigr)$, and the caps cover $\Sph^2$. So
$k\bigl(1-\sqrt{1-\rho^2}\bigr)\ge1$, that is $k\rho^2\ge1+\sqrt{1-\rho^2}\to2$, against $k\rho^2\ge1$
in general. This is the analogue for lines of the cost $\pi$ against $2$ of Proposition~3.
\end{remark}

\subsection{Every convex body: \texorpdfstring{$\lambda_d(n;\Omega)=\Theta(n^{-d})$}{lambda\_d(n;Omega) = Theta(n\textasciicircum -d)} and \texorpdfstring{$\mu_d(d+1;\Omega)=w(\Omega)/(2(d+1))$}{mu\_d(d+1;Omega) = w(Omega)/(2(d+1))}}
\label{sec:D-convex}

This generalizes the order $\lambda_d(n)=\Theta(n^{-d})$ (Corollary~C) and the points-anywhere results
from the unit ball to every convex body. It uses only Bang's plank theorem, Theorem~B's construction with
Lemma~2, and the construction of Section~\ref{sec:D-anywhere}. Theorem~A's strict gain over Bang is not
generalized (its measure is specific to the ball).

\begin{definition}\leavevmode
\begin{itemize}
  \item A convex body $\Omega\subset\R^d$ ($d\ge2$) is a compact convex set with nonempty interior.
  \item $w(\Omega)$ is its minimal width: the least distance between two parallel supporting hyperplanes. It
    is attained, since the width in direction $u$ is continuous on the compact sphere $\Sph^{d-1}$.
  \item $R(\Omega)=\min_{c\in\R^d}\max_{x\in \Omega}\abs{x-c}$ is its circumradius.
  \item For a finite $P\subset\R^d$, $H_P$ is the set of hyperplanes through $d$ affinely independent
    points of $P$, and $\rho_\Omega(P)=\max_{x\in \Omega}\min_{H\in H_P}\dist(x,H)$ ($=\infty$ if $H_P=\emptyset$).
    The max is attained ($\Omega$ is compact).
  \item $\lambda_d(n;\Omega)=\inf\{\rho_\Omega(P):P\subset \Omega,\ |P|=n\}$ and
    $\mu_d(n;\Omega)=\inf\{\rho_\Omega(P):P\subset\R^d,\ |P|=n\}$. For $\Omega=\ball^d$ these are $\lambda_d(n)$ and
    $\mu_d(n)$.
  \item $p_d(n)$, $\sigma_d$ and $N_d(n)$ are as in the main text and Section~\ref{sec:D-anywhere}.
\end{itemize}
\end{definition}

\paragraph{Bang's plank theorem} (Verreault's survey, Bull.\ LMS 58 (2026), e70230, arXiv:2203.05540,
Theorem~2.2, quoted from the arXiv v2 HTML): ``Let $C\subset\R^d$ be a convex body covered by $n$ planks
$P_j$, $j=1,\dots,n$. Then $w(P_1)+\dots+w(P_n)\ge w(C)$.'' There a (closed) plank is ``the closed set
of points between two parallel hyperplanes'', its width is the distance between them, and $w(C)$ is the
minimal width. The survey attributes the theorem to Bang (1950; simplified 1951).

\begin{proposition}\label{D:convex}
Let $\Omega\subset\R^d$ be a convex body, $d\ge2$.
\begin{enumerate}[label=(\roman*)]
  \item For $n\ge d$: $\mu_d(n;\Omega)\ge w(\Omega)/(2\binom nd)$, and so $\lambda_d(n;\Omega)\ge w(\Omega)/(2\binom nd)$.
  \item For $n\ge d+1$: $\lambda_d(n;\Omega)\le R(\Omega)\sin(\pi/(2p_d(n)))$.
  \item For $n\ge d$: $\mu_d(n;\Omega)\le w(\Omega)/(2N_d(n))$. Hence $\mu_d(d+1;\Omega)=w(\Omega)/(2(d+1))$ and
    $d!\,w(\Omega)/2\le\liminf n^d\mu_d(n;\Omega)\le\limsup n^d\mu_d(n;\Omega)\le d^d\,w(\Omega)/2$.
  \item $\lambda_d(n;\Omega)=\Theta(n^{-d})$:
    $d!\,w(\Omega)/2\le\liminf n^d\lambda_d(n;\Omega)\le\limsup n^d\lambda_d(n;\Omega)\le\pi R(\Omega)/(2\sigma_d)$.
\end{enumerate}
\end{proposition}

For $\Omega=\ball^d$ ($w=2$, $R=1$) these are the bounds $1/\binom nd\le\lambda_d(n)\le\sin(\pi/2p_d(n))$,
$\mu_d(d+1)=1/(d+1)$, and so on.

\begin{proof}
\textbf{(i)} Let $P\subset\R^d$ with $|P|=n$ and $\rho=\rho_\Omega(P)<\infty$. Then $\rho>0$, since finitely
many hyperplanes cannot cover $\Omega$, which has nonempty interior. Since the max defining $\rho$ is attained,
every $x\in \Omega$ is within $\rho$ of some $H\in H_P$, so the closed planks $\{x:\dist(x,H)\le\rho\}$,
$H\in H_P$, of width $2\rho$ cover $\Omega$. There are at most $\binom nd$ of them. By Bang's plank theorem,
$\binom nd\cdot2\rho\ge w(\Omega)$. Taking the infimum over $P$ gives the bound for $\mu_d(n;\Omega)$, and
$\lambda_d(n;\Omega)\ge\mu_d(n;\Omega)$ since the infimum is over a smaller set.

\textbf{(ii)} \emph{The circumcentre.} $c\mapsto f(c)=\max_{x\in \Omega}\abs{x-c}$ is $1$-Lipschitz and tends
to $\infty$, so its minimum $R(\Omega)$ is attained at some $c^*$. Then $c^*\in \Omega$: otherwise let $y$ be the
nearest point of $\Omega$ to $c^*$; for every $x\in \Omega$, $\ip{x-y}{c^*-y}\le0$, so
$\abs{x-c^*}^2=\abs{x-y}^2+\abs{y-c^*}^2-2\ip{x-y}{c^*-y}>\abs{x-y}^2$, and $f(y)<f(c^*)$, a
contradiction. Since $\Omega$ is the closure of its interior and $f$ is continuous,
$\inf_{c\in\operatorname{int}\Omega}f(c)=R(\Omega)$. So it suffices to show
$\lambda_d(n;\Omega)\le f(z)\sin(\pi/(2N))$ for every $z\in\operatorname{int}\Omega$, where $N=p_d(n)$.

\emph{The construction.} Fix $z\in\operatorname{int}\Omega$ and $\eta>0$ with $B(z,\eta)\subset \Omega$. Take
$S\subset\Z_N$, $|S|=n$, as in Theorem~B, and, as in its proof, $\theta_j$, $\tau_j=-\cot\theta_j$ and
\[
  q_j(r)=\bigl(\kappa\cdot(\tau_j,\dots,\tau_j^{d-2}),\ r/(i-\tau_j)\bigr)\in\R^{d-2}\times\mathbb C=\R^d
  \qquad(j=1,\dots,n;\ 0<r\le1),
\]
with the $n$ values $\tau_j$ distinct. These $n$ points are distinct (their last components
$r/(i-\tau_j)$ are distinct). For every $d$-subset $F$, the $d$ points
$t_j=\kappa(\tau_j,\dots,\tau_j^{d-2})$, $j\in F$, affinely span $\R^{d-2}$ (Vandermonde), with dependency
$c_j=1/\prod_{k\in F,k\ne j}(\tau_j-\tau_k)$, and $w_F=\sum_{j\in F}c_j/(i-\tau_j)\ne0$. So Lemma~2 for
$j$-flats (Section~\ref{sec:B-lemma2}, with $j=d-1$, $m=2$, $a_j=1/(i-\tau_j)$, $s=r$) applies: the
points $q_j(r)$, $j\in F$, are affinely independent, and their hyperplane $H_F(r)$ satisfies
$\dist(\cdot,H_F(r))\to\dist(\cdot,L_F)$ as $r\to0$, where $L_F=(\R^{d-2}\times\{0\})\oplus\R(0,w_F)$.

The convergence is uniform on every ball $B(0,\mathcal R)$, not only on $\ball^d$: in the proof of
Lemma~2(ii), for $\abs y\le\mathcal R$,
$\abs{(I-\Pi_s)(y-q)-(I-\Pi)(y-q_0)}\le\norm{\Pi_s-\Pi}(\mathcal R+\abs q)+\abs{q-q_0}$, with
$q=q_{i_0}$, $q_0=(t_{i_0},0)$, using $\norm{I-\Pi}\le1$.

By the proof of Theorem~B, the directions of the $w_F$ mod $\pi$, over all $d$-subsets $F$, are $N$
equally spaced ones, so their largest gap mod $\pi$ is $G=\pi/N$.

\emph{Placing it in $\Omega$.} Let $M_0=\sup_{j,\,0<r\le1}\abs{q_j(r)}<\infty$ and $\zeta=\eta/M_0$. The
points $P_j(r)=z+\zeta q_j(r)$ lie in $B(z,\eta)\subset \Omega$, are distinct, and every $d$ of them are
affinely independent, with hyperplane $z+\zeta H_F(r)$. For $x\in \Omega$,
$\dist(x,z+\zeta H_F(r))=\zeta\dist((x-z)/\zeta,H_F(r))$, and $(x-z)/\zeta$ ranges over a bounded set
as $x$ ranges over $\Omega$. So, uniformly in $x\in \Omega$,
$\dist(x,z+\zeta H_F(r))\to\zeta\dist((x-z)/\zeta,L_F)=\dist(x-z,L_F)$, since $L_F$ is a linear
subspace.

\emph{The limit.} A max over $x\in \Omega$ of a min of finitely many functions moves by at most $\eta'$ when
each function moves by at most $\eta'$ uniformly. So
$\rho_\Omega(\{P_j(r)\})\to\max_{x\in \Omega}\min_F\dist(x-z,L_F)$. For $y=x-z$ write
$y=(y_1,y_2)\in\R^{d-2}\times\R^2$. Then $\dist(y,L_F)=\abs{y_2}\sin\angle(y_2,\R w_F)$ ($0$ if
$y_2=0$). Every direction in $\R^2$ is within $G/2=\pi/(2N)\le\pi/2$ of some line $\R w_F$, and $\sin$
increases on $[0,\pi/2]$, so $\min_F\dist(y,L_F)\le\abs{y_2}\sin(\pi/(2N))\le\abs{x-z}\sin(\pi/(2N))$.
Hence $\lambda_d(n;\Omega)\le\lim_{r\to0}\rho_\Omega(\{P_j(r)\})\le f(z)\sin(\pi/(2N))$.

\textbf{(iii)} Choose a unit vector $e$ with width $w(\Omega)$ in direction $e$, and $h_0$ with
$\Omega\subset\{h_0\le\ip xe\le h_0+w(\Omega)\}$. After a rotation and translation, which change neither $\rho_\Omega$
nor the widths, $e=e_d$ and $h_0=0$. Let $N=N_d(n)$ and $S=\{u_1,\dots,u_n\}$ be distinct positive
integers whose $d$-distinct sums contain $c_0+1,\dots,c_0+N$. Apply the construction of
Section~\ref{sec:D-anywhere} (upper bound (i)) with $\alpha=w(\Omega)/N$ and $\beta=-\alpha c_0-\alpha/2$ in
place of $\alpha=2/N$ and $\beta=-1-\alpha c_0-\alpha/2$:
\begin{itemize}
  \item the points $p_i(X)=(Xt_i,\dots,Xt_i^{d-1},\alpha u_i+\beta)$, $t_i=1/u_i$, have every $d$-subset
    affinely independent;
  \item the hyperplane $H_F(X)$ of a $d$-subset $F$ converges in unit normal and offset, as $X\to\infty$,
    to the horizontal hyperplane $\{x_d=\alpha\sum_{i\in F}u_i+\beta\}$. (The proof there works for any
    $\alpha,\beta$: the hyperplane is the graph of $g_F(z)=b_0+\sum_k(b_k/X)z_k$ with
    $b_0=\alpha\sum_{i\in F}u_i+\beta$, and the $b_k$ do not depend on $X$.)
\end{itemize}
The limit list contains the hyperplanes $\{x_d=(j-\frac12)w(\Omega)/N\}$, $j=1,\dots,N$, whose closed
$w(\Omega)/(2N)$-slabs cover $\{0\le x_d\le w(\Omega)\}\supset \Omega$. For $x\in \Omega$, with
$\mathcal R=\max_{x\in \Omega}\abs x$,
\[
  \abs{\dist(x,H)-\dist(x,H')}\le\mathcal R\abs{n-n'}+\abs{o-o'}
\]
for hyperplanes with unit normals $n,n'$ and offsets $o,o'$, and a max-min moves by at most the largest such
change. So $\rho_\Omega(P(X))$ tends to the covering radius of the limit list over $\Omega$, which is at most
$w(\Omega)/(2N)$. Hence $\mu_d(n;\Omega)\le w(\Omega)/(2N_d(n))$.

For $n=d+1$, $N_d(d+1)\ge d+1$ (Section~\ref{sec:D-anywhere}), and (i) gives the matching lower bound
$w(\Omega)/(2\binom{d+1}d)=w(\Omega)/(2(d+1))$. For the limits: $N_d(dq)\ge q^d$ (Section~\ref{sec:D-anywhere},
for $q\ge2$; trivially for $q=1$), and $\mu_d(n;\Omega)$ is non-increasing in $n$ (adding a point keeps every
hyperplane), so for $dq\le n<d(q+1)$, $n^d\mu_d(n;\Omega)\le(d(q+1))^dw(\Omega)/(2q^d)\to d^dw(\Omega)/2$; and
$n^dw(\Omega)/(2\binom nd)\to d!\,w(\Omega)/2$.

\textbf{(iv)} The lower bound is (i) with $n^d/\binom nd\to d!$. For the upper bound, $\sin x\le x$
gives $n^d\lambda_d(n;\Omega)\le R(\Omega)\pi n^d/(2p_d(n))$, and $p_d(n)/n^d\to\sigma_d>0$ (Theorem~B and
Corollary~C), so $\limsup n^d\lambda_d(n;\Omega)\le\pi R(\Omega)/(2\sigma_d)$.
\end{proof}

\begin{remark}
The two constants in (iv) depend on $\Omega$ through $w(\Omega)$ and $R(\Omega)$ separately, and their ratio
$(\pi/(d!\,\sigma_d))R(\Omega)/w(\Omega)$ is not bounded over convex bodies: a $1\times\eps$ rectangle has
$w=\eps$ and $R\approx\frac12$. For such thin bodies the construction in (ii), a pencil through a common
$(d-2)$-flat, is far from the lower bound; points anywhere (iii) do achieve the width scale $w(\Omega)$. For
$\Omega=\ball^d$ Theorem~A improves the lower constant $d!$ to $d!(1+q_d)$; we do not know whether every
convex body admits such a strict gain.
\end{remark}
